\chapter{Reproducibilidad}
\label{ap:reproducibilidad}

Todas las tablas y figuras del informe salen de archivos generados por el código del trabajo. Este apéndice describe
la estructura del proyecto, el entorno y los pasos para reproducir los resultados desde cero.

\section{Estructura del proyecto}

\begin{center}
\small
\begin{tabular}{ll}
\toprule
Ruta & Contenido \\
\midrule
\texttt{datos/bonn/\{Z,O,N,F,S\}}     & Bonn, descargado a mano (no se incluye) \\
\texttt{datos/chbmit/}                & EDF de CHB-MIT, unos 8,8~GB (no se incluyen) \\
\texttt{datos/chbmit\_proc/}          & un \texttt{.npz} por registro, \texttt{crisis.csv}, \texttt{indice.csv}, \texttt{particion.json} \\
\texttt{scripts/descargar\_chbmit.py}  & descarga desde PhysioNet y verificación SHA256 \\
\texttt{scripts/preprocesar\_chbmit.py}& filtrado, ventanas de 4~s, espectrogramas y potencia por banda \\
\texttt{scripts/metricas.py}          & métricas por ventana y por evento, suavizado, umbral de Neyman-Pearson \\
\texttt{scripts/entrenar\_cnn.py}      & CNN y comparativa de optimizadores \\
\texttt{scripts/estilo\_graficos.py}   & paleta y estilo común de todas las figuras \\
\texttt{notebooks/src/*.py}           & fuente de cada notebook (formato \emph{percent} de jupytext) \\
\texttt{notebooks/etapa1\_bonn\_fourier.ipynb}        & Fourier, Welch, filtrado y bandas en Bonn \\
\texttt{notebooks/etapa2\_chbmit.ipynb}               & análisis exploratorio de CHB-MIT y particiones \\
\texttt{notebooks/etapa3\_modelos\_clasicos.ipynb}     & logística, SVM, optimalidad, KKT, LOPO, Neyman-Pearson \\
\texttt{notebooks/etapa4\_cnn\_optimizadores.ipynb}    & análisis de la comparativa y evaluación clínica de la CNN \\
\texttt{resultados/etapa\{1..4\}/}     & figuras, tablas CSV y metadatos de cada etapa \\
\texttt{informe/}                     & este documento \\
\bottomrule
\end{tabular}
\end{center}

\paragraph{Correspondencia con los capítulos.} La Etapa 1 alimenta el Capítulo~\ref{cap:fourier}; la Etapa 2, el
Capítulo~\ref{cap:datos}; la Etapa 3, el Capítulo~\ref{cap:convexos} y las verificaciones numéricas del
Capítulo~\ref{cap:optimizacion}; la Etapa 4, los Capítulos~\ref{cap:cnn} y~\ref{cap:resultados}. Las figuras de
\texttt{informe/figuras/etapaN/} son copias de las de \texttt{resultados/etapaN/}.

\section{Entorno}

\begin{center}
\begin{tabular}{ll}
\toprule
Componente & Versión \\
\midrule
Python        & 3.12.10 \\
NumPy         & 2.5.4 \\
SciPy         & 1.18.1 \\
scikit-learn  & 1.9.1 \\
pandas        & 3.1.1 \\
Matplotlib    & 3.11.2 \\
MNE-Python    & 1.13.2 (lectura de EDF) \\
PyTorch       & 2.15.0, versión para CPU \\
jupytext / nbclient & 1.19.6 / 0.11.0 \\
\midrule
Procesador    & AMD Ryzen 7 5700U, 8 núcleos \\
Memoria       & 21,8~GB \\
GPU           & no \\
Sistema       & Windows, PowerShell \\
\bottomrule
\end{tabular}
\end{center}

Las versiones exactas están en \texttt{requirements.txt}. Los resultados que dependen del azar (particiones de la validación
cruzada, inicialización de la CNN, submuestreo de ventanas normales) usan semillas fijas. PyTorch en CPU es determinístico
para estas operaciones, pero distintas versiones de las bibliotecas de álgebra lineal pueden cambiar los últimos decimales.

\section{Datos}

\begin{itemize}
  \item \textbf{Bonn} \parencite{andrzejak2001}: se descargan los cinco archivos comprimidos (Z, O, N, F, S) de la página del
        grupo de Neurofísica de la Universidad de Bonn y se descomprimen en \texttt{datos/bonn/}, una carpeta por set.
  \item \textbf{CHB-MIT} \parencite{shoeb2009,goldberger2000,guttag2010}: el script descarga los registros de
        \url{https://physionet.org/files/chbmit/1.0.0/} y verifica cada archivo contra el \texttt{SHA256SUMS.txt} oficial.
        La base se distribuye con la licencia Open Data Commons Attribution 1.0.
\end{itemize}

\section{Pasos para reproducir}

Desde la carpeta del proyecto, en PowerShell:
\begin{enumerate}
  \item \textbf{Entorno.}
\begin{verbatim}
python -m venv .venv
.\.venv\Scripts\pip install -r requirements.txt `
    --extra-index-url https://download.pytorch.org/whl/cpu
.\.venv\Scripts\python -m ipykernel install --user --name tp-epilepsia
\end{verbatim}
  \item \textbf{Descarga de CHB-MIT} (unos 8,8~GB; si se corta, se reanuda sin volver a bajar lo verificado):
\begin{verbatim}
.\.venv\Scripts\python scripts\descargar_chbmit.py datos\chbmit `
    chb01 chb02 chb03 chb05 chb08 chb10
\end{verbatim}
  \item \textbf{Preprocesamiento} (con 4 procesos en paralelo):
\begin{verbatim}
.\.venv\Scripts\python scripts\preprocesar_chbmit.py `
    datos\chbmit datos\chbmit_proc 4
\end{verbatim}
  \item \textbf{Notebooks de las Etapas 1 a 3}, desde \texttt{notebooks\textbackslash}; para cada uno:
\begin{verbatim}
..\.venv\Scripts\jupytext --to ipynb src\etapa1_bonn_fourier.py `
    -o etapa1_bonn_fourier.ipynb
..\.venv\Scripts\jupyter execute --inplace `
    --kernel_name=tp-epilepsia etapa1_bonn_fourier.ipynb
\end{verbatim}
        La Etapa 3 guarda los puntajes de los modelos convexos en la partición fija, que la Etapa 4 usa para comparar.
  \item \textbf{Entrenamiento de la CNN}, en tres fases (unos 92~s por corrida en CPU):
\begin{verbatim}
.\.venv\Scripts\python scripts\entrenar_cnn.py . busqueda  # 18 corridas
.\.venv\Scripts\python scripts\entrenar_cnn.py . semillas  # 24 corridas
.\.venv\Scripts\python scripts\entrenar_cnn.py . posthoc   # 20 corridas
\end{verbatim}
        La fase \texttt{semillas} elige la tasa de cada optimizador a partir de las corridas de la búsqueda (mayor AUC de
        validación) y la guarda en \texttt{resultados/etapa4/lr\_elegidas.json}. La fase \texttt{posthoc} corre la tasa
        más alta de cada grilla. Solo son nuevas 20 de sus 30 corridas: para SGD la tasa alta coincide con la elegida, y la
        semilla 0 de los demás ya está en la búsqueda. Cada corrida guarda su historia (JSON) y los puntajes del mejor modelo (NPZ); si el JSON ya
        existe, la corrida se saltea, así que el proceso puede cortarse y reanudarse.
  \item \textbf{Notebook de la Etapa 4}, igual que en el paso 4. Lee las corridas guardadas, arma las tablas y figuras de los
        Capítulos~\ref{cap:cnn} y~\ref{cap:resultados} y hace la evaluación clínica del modelo elegido.
\end{enumerate}

\section{Compilación del informe}

El informe está en \texttt{informe/}: \texttt{main.tex}, un archivo por capítulo en \texttt{capitulos/},
\texttt{referencias.bib} y las figuras en \texttt{figuras/}. Para compilarlo en Overleaf: subir la carpeta comprimida
como proyecto nuevo, elegir el compilador pdf\LaTeX{} (la bibliografía usa \texttt{biblatex} con \texttt{biber}, que
Overleaf ejecuta automáticamente) y compilar.
